\begin{titlepage}
    \begin{center}
        % Judul
        \textbf{LAPORAN AUDIT INSTRUMENTASI}\\[0.0\baselineskip]
        \textbf{CASE-BASED LEARNING (CBL) DAN ELABORATIVE PROJECT PENGUKURAN FISIKA}\\[0.0\baselineskip]
        \textbf{``KETEPATAN DAN KETIDAKPASTIAN DALAM PRAKTIKUM NYATA''}\\[0.0\baselineskip]
        \textbf{PASANGAN EKSPERIMEN P1 (DENSITY) DAN M3 (ENERGY)}\\[0.5\baselineskip]

        % Ruler
        \hrule height 2pt
        \vspace{0.5\baselineskip}

        % Mata Kuliah, Tanggal, dan Kelompok
        {\footnotesize
        \begin{tabular*}{\textwidth}{@{\extracolsep{\fill}} l c r @{}}
            Mata Kuliah: Sensor dan Pengukuran (FIT310) & Tanggal: 17 November 2026 & Kelompok: Group I
        \end{tabular*}
        }
        \vspace{2\baselineskip}

        \begin{figure}[H]
            \centering
            \includegraphics[width=6cm, height=6cm]{Images/logo-unair.png}
        \end{figure}
        \vspace{1\baselineskip}

        \begin{table}[H]
            \centering
            \renewcommand{\arraystretch}{1.05}{
                \begin{tabular}{lc}
                    \multicolumn{2}{c}{\textbf{Disusun Oleh (Group I)}} \\
                    Muhammad Abiyu Khansa & (182241036) \\
                    Farras Atha Pandita & (182241040) \\
                    Hendra Maulana & (182241043)
                \end{tabular}}
        \end{table}

        \textbf{Dosen Pembina :} \\
        \centering
        Winarno, S.Si., M.T. (Pertemuan 1--7) \\
        Yhosep Gita Yhun Yhuwana, S.Si., M.T. (Pertemuan 8--14)

        \vfill
        \begin{center}
            \textbf{PROGRAM STUDI FISIKA \\
                    DEPARTEMEN FISIKA \\
                    FAKULTAS SAINS DAN TEKNOLOGI \\
                    UNIVERSITAS AIRLANGGA \\
                    2026}
        \end{center}
    \end{center}
\end{titlepage}
